\section{Implementation Details}
\label{app:implementation}

\subsection{Numerical verification of the Jacobian}

We unit-test the identity $\partial \zq^{\,\ours} / \partial \ze = s\rmat$
by comparing the autograd Jacobian against the analytic reflection
on a single-token, $d=4$ example. With codebook vector $\qvec =
(1.0, 0.5, -0.2, 0.1)$ and encoder output $\ze = (0.9, 0.45, -0.15,
0.08)$, the agreement between the autograd Jacobian and the
analytic $s\rmat$ is $\|J - s\rmat\|_\infty = 1.19 \times 10^{-7}$
in single precision, which is at the level of FP32 rounding.
Forward-output equality $\zq = \qvec$ holds to within $2 \times
10^{-5}$ across a range of inputs.

\subsection{EMA codebook updates under mixed precision}

Codebook buffers are kept in FP32 regardless of any FP16/BF16
autocast context surrounding the forward pass. When an encoder
feature is produced under autocast it is first cast to FP32 before
updating the codebook's $\ell_2$ moving averages, which avoids
underflow in the cluster-size counts and keeps the codebook
representation consistent across training and evaluation.

\subsection{Hyperparameter table}

\begin{table}[h]
\centering
\small
\begin{tabular}{lr}
\toprule
Hyperparameter & Value \\
\midrule
Encoder base channels              & 64 \\
Encoder / decoder channel mults    & $(1, 2, 2)$ \\
Encoder residual blocks per level  & 2 \\
Latent (z) channels                & 128 \\
Quantizer dim $d$                  & 8 \\
Codebook size $K$                  & 1024 \\
Commitment coefficient $\beta$     & 0.25 \\
EMA decay $\gamma$                 & 0.99 \\
Dead-code threshold                & 1.0 \\
Batch size                         & 1024 \\
Peak learning rate                 & $1.6\times 10^{-3}$ \\
LR schedule                        & cosine to $10^{-6}$ \\
Optimiser                          & AdamW, $\beta=(0.9,0.99)$ \\
Gradient clip                      & 1.0 \\
Training steps                     & 12{,}000 \\
Precision                          & FP32 + TF32 \\
\bottomrule
\end{tabular}
\caption{CIFAR-10 training hyperparameters, shared across all four
quantizers.}
\end{table}

\subsection{Training throughput breakdown}

On a single NVIDIA H100 80GB, at batch size $1024$ and FP32+TF32
precision, a forward-backward-optimiser step of the full $5.40$M
parameter autoencoder takes approximately $165$~ms. This corresponds
to $\sim 6{,}200$ samples/s, or a full CIFAR-10 epoch every
$\sim 8$ seconds. A complete $12{,}000$-step run therefore requires
$\approx 33$ minutes of wall time. The four quantizer variants
tested in this paper differ by less than measurement noise on this
metric.

\subsection{Codebook warm-up}

At initialisation the codebook is populated with a small Gaussian
noise; the first training batch triggers a data-dependent
initialisation that replaces the codebook with random rows from the
first mini-batch of encoder features. This recipe is standard for
VQ-VAEs~\citep{razavi2019vqvae2} and is held identical across
baselines and our method. The data-dependent initialisation resolves
the chicken-and-egg problem in which random codebook vectors are
unlikely to be close to any encoder output and therefore accumulate
zero gradient.

\subsection{Additional case studies}

Additional reconstruction case studies on CIFAR-10 validation images
for each of the four quantizers, along with zoomed error maps, will
accompany the final camera-ready version.
